\documentclass[11pt]{article}
\usepackage[margin=1in]{geometry}
\usepackage{amsmath,amssymb,amsthm}
\usepackage[hidelinks]{hyperref}
\newtheorem{theorem}{Theorem}
\newcommand{\C}{\mathbb C}
\newcommand{\A}{\mathcal A}
\newcommand{\norm}[1]{\left\lVert#1\right\rVert}
\title{A continuous cubic obstruction to the C*-algebraic Smale conjecture}
\author{Research proof and Lean formalization}
\date{October 27, 2026}
\begin{document}
\maketitle
\begin{abstract}
We construct a split cubic in $C([0,1],\C)$ for which no critical element
satisfies the norm form of the C*-algebraic Smale mean-value conjecture.
The derivative also splits, both critical elements exist, and the derivative
at the test element is a nonzero invertible constant. The obstruction is that
connectedness prevents the choice of critical branch from changing between
the two endpoints.
\end{abstract}

Krishna's Conjecture~2.1 in \cite{Krishna} asserts that, for a split polynomial
$P$ of degree $n\geq2$ over a commutative C*-algebra and $P'(z)\ne0$,
some critical element $b$ satisfies
\begin{equation}\label{source}
 \frac{\norm{P(z)-P(b)}}{\norm{z-b}}\leq\norm{P'(z)},
\end{equation}
or the sharper bound with coefficient $(n-1)/n$ on the right.
Here $P'$ is the sum of products obtained by deleting one factor.
We refute already \eqref{source}. The source's degree-two theorem and the
classical scalar Smale conjecture are unaffected.

\begin{theorem}\label{main}
Let $\A=C([0,1],\C)$ with its supremum norm, and put
\[
 y(t)=2t-1,\qquad q(t)=\sqrt{y(t)^2+\tfrac1{16}},\qquad z=y.
\]
For the monic split cubic
\begin{equation}\label{cubic}
 P(w)=w(w-\sqrt3\,q)(w+\sqrt3\,q)=w^3-3q^2w,
\end{equation}
the critical elements are exactly $q$ and $-q$.
The derivative at $z$ is the constant $-3/16$, and neither critical element
satisfies \eqref{source}.
\end{theorem}

\begin{proof}
Expansion of the literal deletion derivative gives
\begin{equation}\label{derivative}
 P'(w)=3(w-q)(w+q),\qquad P'(y)=-\tfrac3{16}\,1.
\end{equation}
Consequently both $q$ and $-q$ are critical. Conversely, if $P'(b)=0$, evaluation
at each $t$ in the scalar field gives $b(t)=q(t)$ or $b(t)=-q(t)$.
Since $q(t)>0$, the continuous function $b/q$ takes values in $\{-1,1\}$.
Connectedness of $[0,1]$ makes it constant. Thus $b=q$ or $b=-q$ globally.
This argument uses scalar evaluation before factor cancellation; it does not
treat $C([0,1],\C)$ as an integral domain.

We have $\norm y\leq1$ and $\norm q\leq5/4$, so for either critical element
\begin{equation}\label{distance}
 \norm{y-b}\leq\tfrac94,\qquad
 \norm{P'(y)}\norm{y-b}\leq\tfrac{27}{64}.
\end{equation}
Write $r=\sqrt{17}/4$. At $t=0$, $y=-1$, $q=r$, and
\[
 P(y)(0)-P(q)(0)=\tfrac{35}{16}+2r^3>2.
\]
At $t=1$, $y=1$, $q=r$, and
\[
 P(y)(1)-P(-q)(1)=-\bigl(\tfrac{35}{16}+2r^3\bigr).
\]
Hence $\norm{P(y)-P(b)}>2$ for every critical element, contradicting
\eqref{distance} if \eqref{source} were true.
The denominator is positive: $y=b$ would imply $P'(y)=0$, contrary to
\eqref{derivative}. Finally the inverse of $P'(y)$ is the constant $-16/3$.
\end{proof}

\paragraph{Scope and formal certificate.}
The counterexample lies in a connected commutative unital C*-algebra.
All three roots in \eqref{cubic} and both distinct critical elements are
continuous. Thus it does not rely on nonexistence of critical elements or a
zero derivative at some test fiber. At individual fibers a suitable scalar
critical point can change sign; a global continuous critical element cannot.

The Lean certificate contains 46 declarations in five modules, including four
reused general continuous-selection lemmas. It uses Lean~4.34.1 and mathlib
revision
\[\texttt{d13f23b723b8a846827a245b89c10fc7d3f11612}.\]
The literal product, deletion derivative, monicity, exact degree, supremum-norm
bounds, branch completeness and explicit inverse are proved.
The closed endpoint
\[\texttt{CStarSmale.source\_counterexample}\]
negates the
universal norm assertion over commutative C*-algebras and all $n\geq2$.
The isolated strict replay passes with only \texttt{propext},
\texttt{Classical.choice} and \texttt{Quot.sound}.

Independent source, prior-art and substantive-eligibility review is pending;
no discovery credit or worldwide priority is claimed. The source's separate
higher-order and dual mean-value conjectures, known public refutations of
those statements, the general selection principle and helper declarations
receive no credit here. This manuscript addresses only the norm Smale
assertion \eqref{source}; an order-strong extension is not a separate outcome.

\begin{thebibliography}{1}
\bibitem{Krishna} K.~Mahesh Krishna,
\emph{C*-algebraic Smale Mean Value Conjecture and Dubinin--Sugawa Dual Mean
Value Conjecture}, arXiv:2206.08154v1 (2022), Conjecture~2.1, pp.~4--5.
\href{https://arxiv.org/abs/2206.08154v1}{arXiv:2206.08154v1}.
\end{thebibliography}
\end{document}
